\documentclass[12pt,a4paper]{article}
\usepackage[utf8]{inputenc}
\usepackage[T1]{fontenc}
\usepackage[polish]{babel}
\usepackage{amsmath, amssymb, amsthm}
\usepackage{geometry}
\usepackage{booktabs}
\usepackage{array}
\usepackage{xcolor}
\usepackage{mdframed}
\usepackage{enumitem}
\usepackage{hyperref}

\geometry{margin=2.5cm}

\definecolor{schemecolor}{RGB}{180, 60, 40}
\definecolor{boxbg}{RGB}{255, 245, 240}
\definecolor{correctgreen}{RGB}{0, 140, 60}
\definecolor{attackred}{RGB}{180, 30, 30}

\newmdenv[
backgroundcolor=boxbg,
linecolor=schemecolor,
linewidth=1.5pt,
roundcorner=5pt,
innertopmargin=10pt,
innerbottommargin=10pt,
innerleftmargin=15pt,
innerrightmargin=15pt
]{schemebox}

\newtheorem*{atak}{Atak}
\newtheorem*{wniosek}{Wniosek}

\title{
	\textbf{Lista nr 2 --- Zadanie 6}\\[0.4em]
}
\author{Kajetan Plewa}
\date{}

\begin{document}
	

\section{LAB14}
\subsection{Troche definicji}

\begin{equation}
	<f_1, ... , f_s> = {\sum_{i=1}^{s}}
\end{equation}


Ideal generowany przez rozmaitosc to wszystkie wielomiany ktore zeruja sie w $a_1, ..., a_n$.

Dla twisted cubic mamy $V(y-x^2, z-x^3)$ to ideal jest tkai sam czyli $I(V)= <y-x^2, z - x^3>$. Dowolne $f \in R[x,y,z]$ wyglada jak : $h_1*(y-x^2)+h_2*(z-x^3)+r(x)$. Reszta jest zero dla naszych rozmaitosci

czyli zeby pokazac cw26 trzeba pokazac ze kazdy wielomian $f \in \kappa[x,y]$ ze kazdy wielomian mozemy rozlozyc jako $f(x,y) = (x-y)*h(x,y)+ g(x)

Wezmy sobie jakis wielomian $x^\alphax^\beta$  -- reszta rozumowania na telefonie 08.05.2026 16:07

Mamy wiec udowodnieone ze tak faktycznie mozna go przedstawic 
Wezmy teraz $f \in I(V(x-y))$ tj.
\end{document}